\documentclass[sigconf, nonacm, review]{acmart}
\renewcommand\footnotetextcopyrightpermission[1]{} 
\settopmatter{printacmref=false}
\pagestyle{plain}
\begin{document}


\title{The title}
\subtitle{Accelerator-Centric Computing Ecosystems -- Programming Assignment Report }

%%
%% The "author" command and its associated commands are used to define
%% the authors and their affiliations.
%% Of note is the shared affiliation of the first two authors, and the
%% "authornote" and "authornotemark" commands
%% used to denote shared contribution to the research.


\author{Klemen Janez Poličar}
\email{k.j.policar@student.vu.nl}
\affiliation{%
  \institution{Computer Science, VU Amsterdam}
  \country{}
}

\author{Author 2}
\email{author2@mail.com}
\affiliation{%
  \institution{Computer Science, VU Amsterdam}
  \country{}
}
\author{Author 3}
\email{author3@mail.com}
\affiliation{%
  \institution{Computer Science, VU Amsterdam}
  \country{}
}


%%
%% The abstract is a short summary of the work to be presented in the
%% article.
\begin{abstract}
  [Optional]  a high-level description of your solution, analysis overview, and main results
\end{abstract}




%%
%% This command processes the author and affiliation and title
%% information and builds the first part of the formatted document.
\maketitle

\section{Introduction}
A brief overview of the problem, motivation for GPU acceleration, and the structure of the remainder of the article. Use one short paragraph for each.\\

\noindent
\textit{Note: you are allowed to deviate from this template as long as the same information is included and it is clear where it is discussed.}

\section{Parallelization}
\textbf{Parallelizable sections and identification.}
We analyzed loop dependencies in the reference code and found that Step 2 (spillage computation) and Step 3 (water update) are grid-wide passes where each iteration operates on one cell with the same arithmetic pattern. These phases dominate runtime and are therefore mapped to CUDA kernels. Step 1 (cloud movement and rainfall) was initially kept on the CPU to keep the first GPU version simple and easy to validate.

\textbf{Workload distribution (geometry).}
We use a 2D CUDA launch where each thread is responsible for one matrix cell $(row,col)$. Blocks are configured as $16\times16$, and the grid dimensions are computed from the terrain size. This mapping follows the natural 2D structure of the terrain and provides regular parallel work across the full domain.

\textbf{Applied optimizations.}
As a first optimization step, we move the expensive per-cell phases to the GPU while preserving the original algorithm. Device buffers for ground and auxiliary arrays are allocated once and reused across iterations. We also use low-overhead reductions via atomics for global metrics (maximum spillage per iteration and boundary water loss), and unroll the fixed 4-neighbor loop to reduce instruction overhead.

\section{Implementation details}
If not already included in the previous section.

\section{Evaluation}

\begin{enumerate}
\item \textit{Experimental setup}: hardware configuration, software details (CUDA version, compiler, etc.), and considered performance metrics
\item \textit{Experiments}: 
\label{experiments} describe the experiments you have conducted to analyze the quality of the performance of your system. This must include at least a speedup analysis (vs. CPU sequential implementation) done with the provided reference datasets. Other experiments can include performance analysis with other dataset sizes, breakdown of execution time, and impact of different optimizations.
\end{enumerate}

\section{Discussion and Conclusion}

Describe the challenges, limitations, and possible optimizations to your solution. We expect you to provide a summary of your findings and (Optional but encouraged) lessons learned and time sheets (report the time it took to conduct each major part of the assignment.


%% The next two lines define the bibliography style to be used, and
%% the bibliography file.
\bibliographystyle{ACM-Reference-Format}
\bibliography{sample-base}
\end{document}